\documentclass[a4paper,12pt]{article}
\usepackage[utf8]{inputenc}
\usepackage[T1]{fontenc}
\usepackage[ngerman]{babel}
\usepackage{geometry}
\usepackage{amsmath}
\geometry{margin=2.5cm}

\begin{document}

\section*{Coupled DEM-LBM Method for the Free-Surface Simulation of
Heterogeneous Suspensions}

\textbf{Autoren:} Alessandro Leonardi, Falk K.\ Wittel, Miller Mendoza,
Hans J.\ Herrmann\\
\textbf{Institution:} Institute for Building Materials, ETH Zürich\\
\textbf{Jahr:} 2015 (arXiv:1509.01052)

% ------------------------------------------------------------------
\subsection*{Zusammenfassung}
% ------------------------------------------------------------------

Die Arbeit präsentiert ein hybrides Simulationsframework, das die Diskrete-Elemente-Methode (DEM) mit der Lattice-Boltzmann-Methode (LBM) koppelt, um heterogene Suspensionen mit freier Oberfläche unter Schwerkrafteinfluss zu simulieren. Das zentrale Anliegen ist die physikalisch konsistente Behandlung der Wechselwirkungen zwischen einer granularen Feststoffphase und einer viskosen Flüssigkeitsphase, wie sie etwa in Frischbeton oder natürlichen Schuttströmen auftreten.

\paragraph{Mehrskalen-Klassifikation.}
Ausgangspunkt der Modellierung ist eine phänomenologische Klassifikation der Körner nach ihrer dominanten Wechselwirkungsphysik mithilfe der dimensionslosen Bagnold-Zahl
\begin{equation}
    \mathrm{Ba} = \frac{\rho_s d_s^2 \lambda_s^{1/2}\,\dot{\gamma}}{\mu_f},
\end{equation}
welche das Verhältnis von kollisiven zu viskosen Kornspannungen beschreibt.
Kleine Körner (elektrostatisch dominiert, $\mathrm{Ba} < 40$) und mittelgroße
Körner (viskos dominiert) werden in der Flüssigphase homogenisiert und durch
ein nicht-Newtonsches Bingham-Kontinuumsmodell mit Fließgrenze $\sigma_y$ und
plastischer Viskosität $\mu_\mathrm{pl}$ beschrieben. Lediglich die
großkörnige Fraktion ($\mathrm{Ba} \geq 450$), bei der kollisive Effekte
dominieren, wird explizit diskret modelliert. Dieses Mehrskalen-Prinzip
ermöglicht eine effiziente Behandlung realistischer Korngrößenverteilungen,
ohne jedes einzelne Partikel auflösen zu müssen.

\paragraph{Diskrete-Elemente-Methode (DEM).}
Die großen Körner werden als Lagrange-Elemente mit translatorischen und
rotatorischen Freiheitsgraden behandelt. Partikelkontakte werden über das
Hertz'sche Kontaktgesetz mit einem Dämpfungsterm modelliert; tangentiale
Kräfte folgen dem Coulomb'schen Reibungsgesetz. Um die für Kugeln typischen
künstlichen Roll- und Mischungsartefakte zu vermeiden, werden zusammengesetzte
nicht-sphärische Elemente aus mehreren Teilkugeln aufgebaut. Die Zeitintegration
erfolgt mit dem Gear-Prädiktor-Korrektor-Schema; die Rotationsdynamik wird
durch Quaternionen singularitätsfrei behandelt.

\paragraph{Lattice-Boltzmann-Methode (LBM).}
Die Flüssigphase wird durch die LBM auf einem D3Q19-Gitter gelöst. Der
kollisive Operator folgt dem BGK-Ansatz (Bhatnagar-Gross-Krook), in dem die
Relaxationszeit $\tau$ direkt mit der lokalen dynamischen Viskosität verknüpft
ist. Die Kopplung der nicht-Newtonschen Bingham-Rheologie erfolgt über eine
ortsabhängige Relaxationszeit $\tau(\mathbf{x}, t)$, die aus dem lokalen
Scherraten-Tensor berechnet wird. Um die Divergenz von $\tau$ für
$\dot{\gamma} \to 0$ zu vermeiden, wird die Viskosität auf einen zulässigen
Bereich $[\tau_\mathrm{min}, \tau_\mathrm{max}]$ begrenzt, was ein
Tri-Viskositäts-Fluid ergibt.

\paragraph{Fluid-Partikel-Kopplung.}
Die no-slip-Randbedingung an beweglichen Partikeloberflächen wird durch die
\emph{Bounce-Back}-Regel realisiert: Eine in Richtung eines Festknotens
strömende Verteilungsfunktion $f_i$ wird in Gegenrichtung reflektiert,
korrigiert um den Wandgeschwindigkeitsterm $6w_i\rho\,\mathbf{u}_w\cdot
\mathbf{c}_i$. Der integrierte Impulsaustausch liefert gleichzeitig die auf
jedes Partikel wirkende Strömungskraft. Bei sehr kleinen Partikelabständen
unterhalb der Gitterknotenauflösung wird zusätzlich eine Lubrikationskraftkorrektur nach Nguyen und Ladd eingesetzt. Fluid- und Festknoten werden bei
jeder Partikelbewegung dynamisch reklassifiziert; neu entstehende Fluidknoten
erhalten ihre Anfangsverteilung durch Mittelung über die Nachbarschaft.

\paragraph{Freie Oberfläche.}
Das Massen-Tracking-Algorithmus von Körner et al.\ unterteilt Fluidknoten in
drei Klassen (flüssig, Interface, Gas). Über eine lokale Massenbilanz der
Verteilungsfunktionen wird die zeitliche Entwicklung der freien Oberfläche
verfolgt; fehlende Verteilungsfunktionen an Interfaceknoten werden aus den
makroskopischen Größen rekonstruiert.

\paragraph{Experimentelle Validierung.}
Das Modell wird anhand eines schwerkraftgetriebenen Ausbreitungsexperiments
mit Frischbeton (Portland-Zement mit 1000 Kieselkieseln, $R = 4$--$8\,\rm
mm$, $\rho_s = 2680\,\rm kg/m^3$) auf einer geneigten Ebene ($15^\circ$)
validiert. Die Simulation auf einem $350 \times 250 \times 80$-LBM-Gitter mit
1000 DEM-Elementen (je 4 Teilsphären) reproduziert die experimentell beobachtete Endform der Betonprobe mit sehr guter Übereinstimmung. Die Gesamtrechenzeit beträgt 63 Stunden auf vier CPU-Kernen -- ein Wert, der die hohen
Kosten vollständig aufgelöster Partikel-Fluid-Simulationen exemplarisch
verdeutlicht.

% ------------------------------------------------------------------
\subsection*{Bezug zur Masterarbeit}
% ------------------------------------------------------------------

Die Arbeit von Leonardi et al.\ ist in der Masterarbeit explizit als
Referenzpunkt für den Rechenaufwand vollständig gekoppelter
Partikel-Fluid-Simulationen zitiert \cite{leonardi_coupled_2014} und
steht in mehrfacher Hinsicht in direktem Bezug zu der dort entwickelten
Methodik.

\paragraph{1. Gemeinsamer physikalischer Kontext: viskose Partikelströmung.}
Sowohl das DEM-LBM-Framework als auch die in der Masterarbeit untersuchten
LaMEM-Simulationen modellieren Partikel, die in einer viskosen Flüssigphase
sedimentieren. Beide Ansätze behandeln das Problem im Regime kleiner
Reynolds-Zahlen (Stokes-Strömung), in dem viskose Kräfte gegenüber Trägheitskräften dominieren und das Strömungsfeld quasistatisch auf Geometrieänderungen reagiert. Das wesentliche physikalische Phänomen -- die Fernwirkung
hydrodynamischer Wechselwirkungen zwischen mehreren Partikeln --
ist in beiden Fällen zentral.

\paragraph{2. Begründung der Surrogat-Modellierung.}
Die von Leonardi et al.\ berichtete Rechenzeit von 63 Stunden für eine
einzelne Simulation mit 1000 Partikeln auf vier CPU-Kernen illustriert
prägnant den Rechenaufwand, dem vollständig aufgelöste Methoden
unterliegen. Die Masterarbeit verweist darauf explizit: LBM-DEM- und
Immersed-Boundary-Ansätze benötigen je Partikel $10^5$--$10^6$ Fluidknoten,
was für Systeme mit mehr als zehn Kristallen zu Gittergrößen von $10^7$--$10^8$
Freiheitsgraden und Laufzeiten von Stunden bis Tagen führt
\cite{leonardi_coupled_2014}. Genau dieser Engpass motiviert die Entwicklung
von Surrogat-Modellen: Neuronale Netze sollen das Strömungsfeld bei
gegebener Kristallkonfiguration in Millisekunden anstatt Stunden
approximieren, ohne die volle Partikel-Fluid-Simulation zu wiederholen.

\paragraph{3. Mehrskalen-Prinzip vs.\ Generalisierungsproblem.}
Die Mehrskalen-Behandlung in Leonardi et al.\ -- feine Partikel werden in
der Flüssigphase homogenisiert, grobe explizit modelliert -- entspricht
konzeptionell der zentralen Fragestellung der Masterarbeit: Kann ein neuronales
Netz, das auf Strömungsfeldern mit $N$ Kristallen trainiert wurde, auf
Konfigurationen mit einer anderen Anzahl $N' \neq N$ generalisieren? Während
Leonardi et al.\ die Mehrskaligkeit durch eine physikalisch begründete
Modelltrennung auflösen, versucht die Masterarbeit, das Netz durch
Trainingsdaten aus einem weiten Bereich von Kristallzahlen
$\{1, \ldots, N_\mathrm{max}\}$ zur impliziten Abstraktion zu befähigen.
Beide Ansätze adressieren somit dasselbe Grundproblem: die Beschreibung
komplexer Partikel-Fluid-Systeme mit variierender Kompositionsstruktur.

\paragraph{4. Behandlung der Randbedingungen und Diskretisierung.}
Die Bounce-Back-Regel der LBM und die Lubrikationskorrektur für kleine
Partikelabstände sind in Leonardi et al.\ notwendig, weil die Partikel
auf dem Eulerschen Gitter diskretisiert werden und ihr Einfluss auf das
Strömungsfeld explizit erzwungen werden muss. In der Masterarbeit übernimmt
LaMEM diese Aufgabe durch scharfe Viskositäts- und Dichtekontraste zwischen
Kristall- und Matrixphase auf einem regulären kartesischen Gitter. Die
konzeptionelle Parallele liegt in der Notwendigkeit, Strömungsgradienten in
unmittelbarer Nähe der Partikeloberflächen korrekt aufzulösen -- eine
Anforderung, die auch für die neuronalen Netze eine besondere Herausforderung
darstellt, da lokale Fehler nahe der Kristallränder in den abgeleiteten
Geschwindigkeitsfeldern $\psi_x, \psi_z$ stärker ins Gewicht fallen.

\paragraph{5. Nicht-newtonsche Rheologie vs.\ newtonsche Grundannahme.}
Das Bingham-Modell in Leonardi et al.\ erlaubt die Simulation von
Fließgrenzenflüssigkeiten, wie sie in Frischbeton oder Schuttströmen vorkommen.
Die Masterarbeit hingegen beschränkt sich bewusst auf einen newtonschen
Matrixfluss konstanter Viskosität $\mu$, um das Generalisierungsproblem über
die Kristallzahl methodisch isoliert untersuchen zu können. Die Erweiterung auf
nicht-newtonsche Matrizes -- etwa zur realistischeren Abbildung von
Magmasystemen mit temperaturabhängiger Viskosität -- bleibt für zukünftige
Arbeiten offen und könnte durch das in Leonardi et al.\ etablierte Bingham-LBM-Konzept informiert werden.

\paragraph{6. Freie Oberfläche: ein abgegrenzter Unterschied.}
Der Massen-Tracking-Algorithmus für freie Oberflächen ist in Leonardi et al.\
wesentlich für die Simulation von Ausbreitungsexperimenten. In der
Masterarbeit werden hingegen stationäre, volumenerhaltende Strömungsfelder
in einem geschlossenen Rechengebiet betrachtet, in dem kein freier Flächenrand
existiert. Dieser Unterschied ist konzeptionell bedeutsam: Die
Stokes-Strömung um sedimentierende Kristalle in der Magmakammer ist
quasistatisch und bildet keine dynamische Phasengrenze aus, wodurch die
Stream-Funktion $\psi$ als globales, divergenzfreies Lernziel formulierbar ist.

\paragraph{Fazit.}
Leonardi et al.\ (2014) liefern sowohl einen methodischen als auch einen
motivationalen Beitrag zur Masterarbeit. Methodisch demonstrieren sie, wie
vollständig aufgelöste Partikel-Fluid-Kopplungen im viskosen Regime realisiert
und validiert werden können -- ein Maßstab, an dem die Qualität der
LaMEM-Trainingsdaten gemessen wird. Motivational verdeutlichen sie durch ihren
Rechenaufwand, warum der Übergang von vollständigen Direktsimulationen zu
physikalisch informierten Surrogat-Modellen in der Geowissenschaft notwendig
und lohnend ist.

\end{document}